\section{问题重述}
\subsection{问题背景}
大语言模型训练需要同时决定参数规模 $N$、训练 token 数 $D$、语料质量 $Q$ 和各领域所占份额 $\mathbf p$。经典标度律描述损失随 $N,D$ 的下降，但无法单独回答质量与配比如何影响损失，以及有限算力应分配给扩模、增加数据、提质或更长上下文。赛题提供附件 A 的质量信号及配比实验、附件 B 的规模与质量实验，以及附件 C 的模型元数据与能力评测；三类数据的来源和可信范围不同，须先统一口径，再逐问连接。

\subsection{问题描述}
\textbf{问题一：}将 A1--A3 的 22 个原始质量字段整理为可比较的指标，构造文档和领域的综合质量分；诊断指标分歧，在 A2/A3 复验筛查规则。结合 A16/A18 将质量信息桥接到 17 个配比领域，以 A4/A5 的训练实验拟合配比对 13 个验证域交叉熵的影响，在 A6--A11 检验并以 A12--A15 的估算数据讨论外推边界，给出候选训练配比。

\textbf{问题二：}以问题一的质量分与配比为接口，建立含参数量、数据量、质量和配比的广义标度律；估计并检验其参数，分析边际效用、领域关系和质量与规模的替代条件，同时明确 A、B 两种质量变量的映射假设。

\textbf{问题三：}在训练、提质和长上下文注意力的联合算力约束下，求不同预算与提质成本函数对应的最优资源配置，并定义和识别配置发生结构性转移的临界预算。

\textbf{问题四：}基于附件 C 建立损失与能力评测的桥接，区分规模扩张与非规模技术进步的贡献，在算力增长放缓情景下预测开源模型能力前沿并说明不确定性。

\section{问题分析}
\subsection{总体思路}
四问按接口递进：问题一提供唯一主质量分 $Q$、17 域质量分和配比—损失关系；问题二在可识别范围内构造广义标度律；问题三将该损失函数置于预算约束下优化；问题四把预测损失映射为能力分数。A1 拟合的预处理、权重、Huber 参数和冲突筛查参照冻结用于 A2/A3；A4/A5 用于配比拟合，A6/A7 作 1M 留出；A8/A9 的 60M 数据用于冻结检验及尺度修正标定，A10/A11 的 1B 数据留作独立检验。半合成数据用于模型形式比较，估算表不得视为真实大模型训练验证。

\subsection{各问题的关键难点}
\textbf{问题一：}质量指标量纲、方向和适宜区间不同，DSIR 与文档长度高度相关；同一来源内部也可能意见分裂，不能用一个来源中心替代全部指标。17 个配比域中仅 6 个有直接或近似的质量映射；份额之和为 1，直接回归会忽略成分约束。分别采用 A1 冻结预处理、CRITIC--Huber 主分、评分后的单指标与六来源分层诊断、文本规则特征桥接及 clr 变换。

\textbf{问题二：}B 组数据性质不一致，A、B 的质量变量没有可直接估计的共同锚点；质量项形式和跨来源适用性必须分开检验。\textbf{问题三：}成本约束非线性且质量边界可能成为最优点，需要同时检查解析条件与数值解。\textbf{问题四：}能力分数与训练损失口径不同，时间、规模及模型类型相互混杂，预测应按可比性与情景分层。
